\documentclass[11pt,a4paper]{article}
\usepackage[utf8]{inputenc}
\usepackage[T1]{fontenc}
\usepackage[polish]{babel}
\usepackage{lmodern}
\usepackage[margin=27mm]{geometry}
\usepackage{microtype}
\usepackage[hidelinks]{hyperref}
\title{Błędy predykcji społecznej i asymetria niedopasowania wyglądu i zachowania}
\author{Mateusz Pomianek}
\date{Szkic naukowy, wersja 0.1.1, 21 listopada 2026}
\begin{document}
\maketitle
\begin{abstract}
Czy ten sam klip docelowy wywołuje różne oceny po różnym kontekście oczekiwań i czy efekt zależy od kierunku niezgodności wyglądu z ruchem? Przedstawiono prospektywny projekt badania oparty na wspólnej platformie ST-08, wykorzystujący archiwalne nagrania zachowań odtwarzane na syntetycznych postaciach o kontrolowanym wyglądzie. Pierwszorzędowy wskaźnik: Interakcja kontekstu oczekiwań oraz kierunku niedopasowania w ocenie niesamowitości klipu docelowego. Badanie nie dostarczyło dotąd danych ani wyników empirycznych.
\end{abstract}
\section{Wprowadzenie}
Zależność między wyglądem jednostki a widocznym przebiegiem jej zachowania jest zasadniczym elementem społecznej percepcji. W kontaktach z robotami społecznymi i realistycznymi awatarami nieoczekiwana mimika, gest, kontakt wzrokowy i dynamika wypowiedzi mogą wpływać na reakcje odbiorcy. Nie oznacza to jednak, że każda atypowa ekspresja musi być oceniana negatywnie, ani że jest uprawnione utożsamianie ekspresji z intensywnością doświadczanych emocji.

Projekt stanowiska ST-08 umożliwi oddzielenie zachowania rzeczywistego człowieka, wyodrębnionego z istniejących nagrań interakcji, od syntetycznej formy jego reprezentacji. Pozwala to porównywać te same zdarzenia w kontrolowanym wyglądzie awatara. Nie można jednak przyjąć, że retargeting automatycznie zachowuje wszystkie własności zachowania. Jego dokładność i sposób deformowania sygnałów muszą być sprawdzone niezależnie od właściwego eksperymentu.
\section{Cel i hipotezy}
Czy ten sam klip docelowy wywołuje różne oceny po różnym kontekście oczekiwań i czy efekt zależy od kierunku niezgodności wyglądu z ruchem?
\paragraph{H1.} Manipulacja oczekiwaniem społecznym zmieni ocenę niesamowitości dla identycznego docelowego zachowania.
\paragraph{H2.} Niedopasowanie robotyczny wygląd z ludzkim ruchem nie będzie równoważne ludzkiemu wyglądowi z ruchem mechanicznym.
Przed badaniem potwierdzającym konieczne jest ustalenie estymandu, sposobu kodowania efektu, kryteriów jakości, mocy oraz procedury radzenia sobie z wielokrotnymi porównaniami.
\section{Materiał i metody}
Identyczne klipy docelowe na różnych awatarach, poprzedzone kontekstami przewidującymi odmienne reakcje. Utrzymywana jest ta sama treść i czas klipu docelowego.

Schemat czynnikowy kierunek niedopasowania razy zgodność z uprzednio wywołanym oczekiwaniem, z kontrolą powtórzeń, ekspozycji i nowości.

Pierwszorzędowy wskaźnik: Interakcja kontekstu oczekiwań oraz kierunku niedopasowania w ocenie niesamowitości klipu docelowego. Jednostką analizy jest ocena konkretnej sekwencji przez konkretnego obserwatora; należy uwzględnić zależność prób od źródłowej osoby, scenariusza i postaci. Zbieranie odpowiedzi behawioralnych nie wymaga EEG ani okulografii. Szacowana orientacja spojrzenia w archiwalnym filmie nie jest laboratoryjnym pomiarem eye trackingu.
\section{Plan analizy}
Model mieszany różnic różnic z efektami uczestnika, klipu i kontekstu; osobno zweryfikować skuteczność manipulacji i kontrolować alternatywne wyjaśnienia.

Zamrożenie estymandu i planu porównań nastąpi po pilotażu technicznym i przed dostępem do danych konfirmacyjnych. Należy zachować oddzielnie analizy eksploracyjne, raportować przedziały niepewności, wielkości efektów oraz reguły wyłączenia bodźców z błędną animacją. Trafność konstruktu wymaga niezależnej kalibracji realizmu wyglądu i percepcyjnej naturalności zachowania.
\section{Wyniki}
Nie zebrano jeszcze danych empirycznych i nie wykonano analiz potwierdzających. W niniejszym opracowaniu nie przypisuje się wartości statystycznych, wyników skal, skuteczności diagnostycznej ani dowodów występowania doliny niesamowitości.
\section{Dyskusja i ograniczenia}
Praca testuje mechanizm predykcyjny, a nie tylko opisuje nieliniową powierzchnię ocen P02. Prawa do korzystania z oryginalnych nagrań i rozpowszechniania pochodnych animacji nie wynikają automatycznie z publicznej dostępności bazy. Rekonstrukcja mimiki i ruchu z pojedynczej kamery może zawierać systematyczne błędy. Dla badań z udziałem nowych obserwatorów konieczne są właściwa kwalifikacja etyczna, ochrona danych i jasno określona procedura zgody.

Wyniki negatywne oraz niejednoznaczne mają znaczenie naukowe, jeśli kontrolowano jakość bodźców, precyzję analizy i heterogeniczność kontekstu. Szczególnie w przypadku różnorodności komunikacji nie można wywodzić wewnętrznych stanów psychicznych z nietypowej ekspresji widocznej na twarzy ani sprowadzać zróżnicowania do jednolitej cechy klinicznej.
\section{Wnioski}
Zaproponowano odrębny problem naukowy i sposób jego przyszłego falsyfikowania. Uprawnienia do nagrań, walidacja rekonstrukcji zachowania, prerejestracja i badania obserwatorów pozostają do wykonania. Powstanie szkicu LaTeX nie stanowi dowodu gotowości publikacyjnej.
\end{document}
